% Vietnamese translation for OI Public Library
% Source-PDF: IOI中国国家候选队论文/2009/方展鹏/浅谈如何解决不平等博弈问题.pdf
\documentclass[11pt]{article}
\usepackage{oipl-vn}

\newcommand{\Gpos}[2]{\left\{\,#1\,\middle|\,#2\,\right\}}
\newcommand{\Oh}{\mathcal{O}}
\newcommand{\Inf}{\mathit{infinite}}
\newcommand{\sgn}{\operatorname{sign}}

\title{Bàn về cách giải bài toán trò chơi bất đối xứng}
\author{Fang Zhanpeng\\Trường Trung học số 1 Trung Sơn, Quảng Đông}
\date{2009}

\begin{document}
\maketitle

\origtitle{Bàn về cách giải bài toán trò chơi bất đối xứng}
\sourcepath{IOI China candidate team papers / 2009 / Fang Zhanpeng}

\begin{abstract}
Bài viết giới thiệu cách xử lý một lớp trò chơi trong đó hai người chơi không
có cùng tập nước đi hợp lệ. Nội dung gồm bốn phần: mở đầu bằng sự khác nhau
giữa các bài toán trò chơi; dùng surreal number để giải một lớp trò chơi tổ
hợp bất đối xứng; quay lại phân tích trạng thái, truy hồi và lặp để xử lý các
bài toán không phù hợp với công cụ đó; cuối cùng là tổng kết.
\end{abstract}

\paragraph{Từ khóa.}
Trò chơi bất đối xứng; surreal number; phân tích trạng thái; truy hồi; lặp.

\section{Mở đầu}

Trong thi Tin học, bài toán trò chơi xuất hiện khá thường xuyên. Xét ví dụ
kinh điển \emph{Green Hackenbush}: cho \(n\) cây tre có chiều cao lần lượt là
\(a_1,a_2,\ldots,a_n\). Hai người chơi \(L\) và \(R\) chơi trò chặt tre theo
các luật sau:
\begin{enumerate}
  \item Hai người lần lượt thao tác, \(L\) đi trước.
  \item Mỗi lượt chọn một cây tre có chiều cao khác \(0\), chặt bỏ một đoạn
  của cây đó; mọi phần không còn nối với gốc cũng bị bỏ đi.
  \item Người đầu tiên không thể thao tác sẽ thua.
\end{enumerate}
Nếu cả hai đều tối ưu, cần xác định ai thắng ở trạng thái đã cho.

Một cây tre cao \(n\) sau một lần chặt có thể trở thành một cây cao từ \(0\)
đến \(n-1\). Vì vậy nó tương đương một đống Nim có \(n\) viên. Bài toán trở
thành Nim với \(n\) đống có kích thước \(a_1,a_2,\ldots,a_n\): mỗi lượt chọn
một đống và lấy đi ít nhất một viên, người không đi được thua. Theo định lý
Sprague--Grundy, nếu \(g_i\) là hàm SG của trò chơi con \(G_i\), thì trò chơi
tổng \(G=G_1+G_2+\cdots+G_n\) có hàm SG
\[
  g(x_1,x_2,\ldots,x_n)
  =g_1(x_1)\operatorname{xor}g_2(x_2)\operatorname{xor}\cdots
   \operatorname{xor}g_n(x_n).
\]
Do đó bài toán trên có thuật toán \(\Oh(n)\).

Trong Green Hackenbush, ở mọi trạng thái, \(L\) và \(R\) có cùng tập lựa chọn:
nếu \(L\) có thể chặt đoạn thứ \(j\) của cây thứ \(i\), thì \(R\) cũng có thể.
Nếu điều này không còn đúng thì sao? Ta thêm một luật: mỗi đoạn tre được gán
nhãn \(L\) hoặc \(R\); người chơi \(L\) chỉ được chặt đoạn mang nhãn \(L\), và
\(R\) chỉ được chặt đoạn mang nhãn \(R\). Khi ấy tập quyết định của hai bên
khác nhau, đồng thời định lý Sprague--Grundy không còn áp dụng trực tiếp.

Bài viết này bàn về cách xử lý những trò chơi trong đó hai người chơi có tập
nước đi hợp lệ khác nhau; ta gọi chúng là trò chơi bất đối xứng
(\emph{partizan games}).

\section{Một hướng khác}

Phần này giới thiệu công cụ \emph{surreal number}, rồi dùng một ví dụ để phân
tích cách giải một lớp trò chơi tổ hợp bất đối xứng.

\subsection{Giới thiệu surreal number}

Trong toán học, surreal number tạo thành một trường chứa cả số thực, số vô
cùng lớn và số vô cùng bé. Ở đây ta chỉ dùng nó như công cụ phân tích trò
chơi, nên các tính chất ít liên quan sẽ không được trình bày.

\subsubsection{Cách xây dựng}

\paragraph{Định nghĩa 1.}
Một surreal number gồm hai tập, gọi là tập trái và tập phải. Ta thường viết
nó dưới dạng \(\Gpos{L}{R}\), trong đó \(L\) là tập trái và \(R\) là tập phải.
Các phần tử trong hai tập cũng là surreal number, đồng thời không tồn tại
\(x\in R\) và \(y\in L\) sao cho \(x\le y\).

\paragraph{Định nghĩa 2.}
Với hai surreal number
\[
  x=\Gpos{X_L}{X_R},\qquad y=\Gpos{Y_L}{Y_R},
\]
ta nói \(x\le y\) khi và chỉ khi không tồn tại \(x_L\in X_L\) sao cho
\(y\le x_L\), và cũng không tồn tại \(y_R\in Y_R\) sao cho \(y_R\le x\).
Để tiện trình bày, viết \(x\ge y\) thay cho \(y\le x\), viết
\(x<y\) thay cho \(x\le y\land \neg(y\le x)\), viết \(x>y\) thay cho \(y<x\),
và viết \(x=y\) thay cho \(x\le y\land y\le x\).

Hai định nghĩa trên đều mang tính đệ quy. Ban đầu ta chưa biết surreal number
nào, nên chỉ có thể lấy cả hai tập rỗng, thu được \(\Gpos{\varnothing}
{\varnothing}\), thường viết gọn là \(\Gpos{}{}\). Theo định nghĩa 1, đây là
một surreal number hợp lệ; đặt \(\Gpos{}{}=0\), và nói \(0\) sinh ở ngày thứ
\(0\).

Từ \(0\), ta thử tạo ra các số mới:
\[
  \Gpos{0}{},\qquad \Gpos{}{0},\qquad \Gpos{0}{0}.
\]
Vì \(0\le 0\), \(\Gpos{0}{0}\) không hợp lệ. Theo định nghĩa 2,
\[
  \Gpos{}{0}<0<\Gpos{0}{},
\]
nên đặt \(\Gpos{}{0}=-1\), \(\Gpos{0}{}=1\). Hai số \(1\) và \(-1\) sinh ở
ngày thứ \(1\).

Dùng \(0,1,-1\) làm phần tử của tập trái và tập phải, ta dựng được \(17\)
surreal number hợp lệ ở ngày thứ \(2\). Chẳng hạn
\[
  \Gpos{1}{}=2,\qquad \Gpos{}{-1}=-2.
\]
Ngoài ra \(0<\Gpos{0}{1}<1\). Vì phép cộng sẽ cho
\(\Gpos{0}{1}+\Gpos{0}{1}=1\), ta đặt \(\Gpos{0}{1}=1/2\); tương tự
\(\Gpos{-1}{0}=-1/2\). Như vậy ta đã có
\[
  -2,\ -1,\ -\frac12,\ 0,\ \frac12,\ 1,\ 2.
\]
Những số hợp lệ còn lại ở ngày thứ \(2\) có giá trị trùng với một trong các
số trên; ví dụ \(\Gpos{-1}{1}=0\).

Tiếp tục như vậy, ngày thứ \(3\) cho
\[
  \frac14=\Gpos{0}{1/2},\quad
  \frac34=\Gpos{1/2}{1},\quad
  \frac32=\Gpos{1}{2},\quad
  3=\Gpos{2}{},
\]
cùng các số âm tương ứng. Lặp mãi, ta xây dựng được mọi hữu tỉ dạng
\(j/2^k\). Có thể mô tả một phần ánh xạ từ các hữu tỉ này sang surreal number
bằng hàm dyadic \(\delta\):
\[
\delta(x)=
\begin{cases}
  \Gpos{}{}, & x=0,\\
  \Gpos{\delta(x-1)}{}, & x>0,\ x\in\mathbb{Z},\\
  \Gpos{}{\delta(x+1)}, & x<0,\ x\in\mathbb{Z},\\
  \Gpos{\delta\!\left(\frac{j-1}{2^k}\right)}
        {\delta\!\left(\frac{j+1}{2^k}\right)}, &
  x=\frac{j}{2^k},\ j,k\in\mathbb{Z},\ k>0.
\end{cases}
\]
Thực ra surreal number còn dựng được toàn bộ số thực, số vô cùng lớn và số vô
cùng bé, nhưng phần đó nằm ngoài phạm vi bài viết.

\subsubsection{Một số định lý cơ bản}

\paragraph{Định lý 1.}
Với surreal number \(x=\Gpos{L}{R}\), \(x\) lớn hơn mọi phần tử của \(L\) và
nhỏ hơn mọi phần tử của \(R\).

\paragraph{Định lý 2.}
Với \(x=\Gpos{L}{R}\), nếu \(L\) có phần tử lớn nhất \(l_{\max}\) thì
\(\Gpos{l_{\max}}{R}=x\). Tương tự, nếu \(R\) có phần tử nhỏ nhất
\(r_{\min}\) thì \(\Gpos{L}{r_{\min}}=x\). Ví dụ:
\[
  \Gpos{1,2,3}{4,5}=\Gpos{3}{4,5}=\Gpos{1,2,3}{4}=\Gpos{3}{4}.
\]
Ta không thể sửa định nghĩa surreal number thành ``mỗi tập trái và phải có
nhiều nhất một phần tử'', vì khi các tập là vô hạn, chúng có thể không có phần
tử lớn nhất hoặc nhỏ nhất.

\paragraph{Định lý 3.}
Giả sử surreal number sinh sớm nhất trong khoảng từ \(a\) đến \(b\) sinh ở
ngày thứ \(i\). Khi đó trong khoảng ấy chỉ có đúng một surreal number sinh ở
ngày thứ \(i\).

\paragraph{Định lý 4.}
Nếu \(a<x<b\), và \(x\) là surreal number sinh sớm nhất giữa \(a\) và \(b\),
thì \(\Gpos{a}{b}=x\).

\subsubsection{Phép toán}

Các phép cộng và trừ của surreal number cũng được định nghĩa đệ quy.

\paragraph{Định nghĩa 3.}
Với
\[
  x=\Gpos{X_L}{X_R},\qquad y=\Gpos{Y_L}{Y_R},
\]
đặt
\[
  x+y
  =\Gpos{X_L}{X_R}+\Gpos{Y_L}{Y_R}
  =\Gpos{X_L+y,\ x+Y_L}{X_R+y,\ x+Y_R},
\]
trong đó với tập \(X\) và surreal number \(y\),
\[
  X+y=\{x+y:x\in X\}.
\]
Điều kiện dừng của đệ quy là \(\varnothing+n=\varnothing\).

\paragraph{Định nghĩa 4.}
Với \(x=\Gpos{X_L}{X_R}\), số đối của \(x\) là
\[
  -x=-\Gpos{X_L}{X_R}=\Gpos{-X_R}{-X_L},
\]
trong đó \(-X=\{-x:x\in X\}\). Điều kiện dừng là
\[
  -0=-\Gpos{}{}=\Gpos{}{}=0.
\]

\paragraph{Định nghĩa 5.}
Với surreal number \(x,y\), đặt \(x-y=x+(-y)\).

Từ các định nghĩa trên suy ra:
\begin{enumerate}
  \item Nếu \(x=\Gpos{X_L}{X_R}\) và \(y=\Gpos{Y_L}{Y_R}\), thì
  \[
    x+y=\Gpos{X_L+y,\ x+Y_L}{X_R+y,\ x+Y_R}
  \]
  là một surreal number hợp lệ.
  \item Phép cộng giao hoán: \(x+y=y+x\).
  \item Phép cộng kết hợp: \((x+y)+z=x+(y+z)\).
\end{enumerate}

\subsection{Surreal number và trò chơi tổ hợp}

\subsubsection{Định nghĩa và biểu diễn trò chơi tổ hợp}

Trước khi áp dụng surreal number, ta cố định lớp trò chơi được xét:
\begin{itemize}
  \item Có hai người chơi, gọi là \(L\) và \(R\).
  \item Ở mọi thời điểm, trò chơi có một trạng thái xác định.
  \item Một thao tác chuyển trò chơi từ trạng thái hiện tại sang trạng thái
  khác; luật chơi cho biết từ mỗi trạng thái người chơi có thể tới những trạng
  thái nào.
  \item Hai người chơi luân phiên thao tác.
  \item Nếu tới lượt một người mà người đó không thể thao tác, trò chơi kết
  thúc. Ở đây chỉ xét quy ước người đầu tiên không thao tác được sẽ thua.
  \item Với mọi chuỗi thao tác, trò chơi luôn kết thúc sau hữu hạn bước, tức
  không có hòa.
  \item Hai người chơi biết đầy đủ luật chơi, trạng thái và mọi thao tác trước
  đó.
\end{itemize}

Vì trò chơi được tạo bởi các trạng thái, nếu ở trạng thái \(P\), tập trạng thái
mà \(L\) có thể chuyển tới là \(P_L\), còn tập trạng thái mà \(R\) có thể
chuyển tới là \(P_R\), ta viết
\[
  P=\Gpos{P_L}{P_R}.
\]
Ví dụ, nếu từ \(P\), \(L\) có thể tới \(A,B,C\), còn \(R\) có thể tới \(D\),
thì
\[
  P=\Gpos{A,B,C}{D}.
\]

Cách viết này rất giống surreal number. Thật vậy, mọi surreal number đều là
một trò chơi, nhưng không phải trò chơi nào cũng là surreal number, vì với
trò chơi \(P=\Gpos{P_L}{P_R}\) không bắt buộc mọi phần tử của \(P_L\) nhỏ hơn
mọi phần tử của \(P_R\). Sau đây ta chỉ xét các trò chơi có thể biểu diễn bằng
surreal number.

\subsubsection{Áp dụng vào trò chơi tổ hợp}

Giả sử trò chơi \(G\) tương đương surreal number \(x\), và hai người chơi đều
tối ưu. Khi đó:
\begin{itemize}
  \item Nếu \(x>0\), người chơi \(L\) thắng dù đi trước hay đi sau.
  \item Nếu \(x<0\), người chơi \(R\) thắng dù đi trước hay đi sau.
  \item Nếu \(x=0\), người đi sau thắng.
\end{itemize}

Các kết luận này chứng minh được bằng ý tưởng quy nạp. Khi \(x>0\), nếu \(L\)
đi trước thì vì \(x\) là surreal number, \(L\) luôn có thể chuyển tới một
trạng thái \(\ge 0\). Nếu \(R\) đi trước, \(R\) chỉ có thể chuyển tới một
trạng thái \(>0\). Vì vậy tập nước đi của \(L\) không bao giờ rỗng trước, và
\(L\) thắng. Trường hợp \(x<0\) đối xứng. Khi \(x=0\), nếu \(L\) đi trước thì
chỉ có thể chuyển sang trạng thái \(<0\), nên \(R\) thắng; nếu \(R\) đi trước
thì chỉ có thể chuyển sang trạng thái \(>0\), nên \(L\) thắng.

Nhiều trò chơi \(G\) có thể tách thành \(n\) trò chơi con rời nhau
\(G_1,G_2,\ldots,G_n\), mỗi thao tác của \(G\) tương đương chọn một trò chơi
con rồi thao tác trong đó. Khi ấy ta gọi \(G\) là tổng của các trò chơi con,
viết
\[
  G=G_1+G_2+\cdots+G_n.
\]
Nim nhiều đống là ví dụ kinh điển.

\paragraph{Định lý 8.}
Nếu trò chơi \(G\) tương đương surreal number \(x\), và trò chơi \(H\) tương
đương surreal number \(y\), thì trò chơi \(G+H\) tương đương surreal number
\(x+y\).

Định lý này đến trực tiếp từ định nghĩa phép cộng. Nhờ đó, surreal number giải
được nhiều trò chơi bất đối xứng mà định lý Sprague--Grundy không xử lý được.

\subsection{Ví dụ 1: Procrastination}

\subsubsection*{Tóm tắt đề bài}

Trò chơi \emph{Procrastination} có luật như sau:
\begin{itemize}
  \item Ban đầu có bốn tháp, mỗi tháp là một chồng lập phương đen hoặc trắng.
  \item Hai người chơi \(L\) và \(R\) luân phiên thao tác.
  \item Mỗi lượt, người chơi chọn một lập phương rồi bỏ lập phương đó cùng mọi
  lập phương nằm phía trên nó. Người chơi \(L\) chỉ được chọn lập phương
  trắng; \(R\) chỉ được chọn lập phương đen.
  \item Người đầu tiên không thể thao tác sẽ thua.
\end{itemize}

Một trạng thái ban đầu được gọi là \emph{W-configuration} nếu \(L\) thắng dù đi
trước hay đi sau. Một cấu hình con \(C\) gồm ba tháp. Một trạng thái đầy đủ có
thể được tạo bởi \(C\) và một tháp thứ tư \(T\), ký hiệu \((C,T)\). Với hai cấu
hình con \(C_1,C_2\), nói \(C_1\) không kém \(C_2\) nếu với mọi tháp \(T\),
hễ \((C_2,T)\) là W-configuration thì \((C_1,T)\) cũng là W-configuration.

Cho \(C_1,C_2\), cần kiểm tra \(C_1\) có không kém \(C_2\) hay không. Mỗi tháp
đầu vào có chiều cao \(n\), với \(0\le n\le 50\).

\subsubsection*{Phân tích bằng surreal number}

Trước hết xét một trạng thái \(G\) chỉ gồm một tháp \(T\). Nếu \(T\) rỗng thì
cả hai người đều không có nước đi, nên
\[
  G=\Gpos{}{}=0.
\]
Nếu \(T\) chỉ có một lập phương trắng, \(L\) có một nước đi tới \(0\), nên
\[
  G=\Gpos{0}{}=1.
\]
Nếu \(T\) gồm hai lập phương trắng, \(L\) có thể đưa trò chơi tới \(1\) hoặc
\(0\), nên
\[
  G=\Gpos{0,1}{}=2.
\]
Quy nạp cho thấy tháp gồm \(n\) lập phương trắng có giá trị \(n\); tháp gồm
\(n\) lập phương đen có giá trị \(-n\).

Tiếp theo xét tháp gồm \(n>0\) lập phương trắng ở dưới và một lập phương đen ở
trên. Khi đó
\[
  G=\Gpos{0,1,\ldots,n-1}{n}
   =\Gpos{n-1}{n}
   =n-\frac12.
\]
Nếu thêm một lập phương đen nữa lên trên, ta có
\[
  G=\Gpos{n-1}{n-\frac12}
   =n-\frac12-\frac14.
\]
Nếu từ trạng thái đó thêm một lập phương đen ở đỉnh, giá trị thành
\[
  n-\frac12-\frac14-\frac18;
\]
còn nếu thêm một lập phương trắng ở đỉnh, giá trị thành
\[
  n-\frac12-\frac14+\frac18.
\]

Với một tháp bất kỳ, giá trị tương ứng được tính bằng thuật toán sau. Ký hiệu
\(T[i]\) là màu của lập phương thứ \(i\) tính từ dưới lên.

\begin{verbatim}
SurrealNumber(T)
    x <- 0
    i <- 1
    n <- chieu cao cua T
    while i <= n and T[i] = T[1]
        if T[i] = trang then x <- x + 1 else x <- x - 1
        i <- i + 1
    k <- 2
    while i <= n
        if T[i] = trang then x <- x + 1 / k else x <- x - 1 / k
        i <- i + 1
        k <- k * 2
    return x
\end{verbatim}

Thuật toán có hai bước: trước hết xử lý đoạn đáy gồm các lập phương cùng màu,
sau đó xử lý từng lập phương phía trên đoạn đó bằng các trọng số
\(1/2,1/4,1/8,\ldots\). Độ phức tạp là \(\Oh(n)\). Có thể hiểu cấu hình này
như một đoạn đáy cùng màu, rồi các lập phương phía trên lần lượt đóng góp các
lũy thừa âm của \(2\) với dấu phụ thuộc màu.

Chứng minh đúng đắn của bước thứ hai dùng quy nạp. Xét tháp \(T\) có \(n\) lập
phương dưới cùng cùng màu \(C_1\), lập phương thứ \(n+1\) có màu khác \(C_1\),
và phía trên đoạn đáy có tổng cộng \(p\) lập phương.

Với \(p=1\), giả sử đoạn \(T[1..n]\) có giá trị \(v\). Nếu lập phương trên
cùng là đen, thì
\[
  G=\Gpos{0,1,\ldots,n-1}{n}
   =\Gpos{n-1}{n}
   =n-\frac12
   =v-\frac12.
\]
Nếu lập phương trên cùng là trắng, lập luận đối xứng cho
\[
  G=-n+\frac12=v+\frac12.
\]

Giả sử thuật toán đúng với mọi \(p=1,\ldots,m-1\), ta chứng minh cho
\(p=m\). Đánh số \(m+1\) lập phương trên cùng từ trên xuống dưới là
\(m,m-1,\ldots,0\).

Nếu lập phương số \(m\) là đen, tìm từ trên xuống lập phương trắng đầu tiên và
gọi số của nó là \(k\). Gọi \(v\) là giá trị của \(T[1..n+k]\), và \(u\) là giá
trị của \(T[1..n+m-1]\). Theo định lý 2,
\[
\begin{aligned}
G
&=\Gpos{v-\frac1{2^k}}
        {v-\frac1{2^{k+1}}-\frac1{2^{k+2}}-\cdots-\frac1{2^{m-1}}}\\
&=\Gpos{v-\frac1{2^k}}
        {v-\frac1{2^k}+\frac1{2^{m-1}}}\\
&=\Gpos{\dfrac{v2^m-2^{m-k}}{2^m}}
        {\dfrac{v2^m-2^{m-k}+2}{2^m}}.
\end{aligned}
\]
Theo hàm dyadic, số sinh sớm nhất giữa hai cận trên là
\[
  G=\frac{v2^m-2^{m-k}+1}{2^m}
   =v-\frac1{2^k}+\frac1{2^{m-1}}-\frac1{2^m}
   =u-\frac1{2^m}.
\]

Nếu lập phương số \(m\) là trắng, tìm lập phương đen đầu tiên từ trên xuống và
gọi số của nó là \(k\). Khi đó
\[
\begin{aligned}
G
&=\Gpos{v+\frac1{2^{k+1}}+\frac1{2^{k+2}}+\cdots+\frac1{2^{m-1}}}
        {v+\frac1{2^k}}\\
&=\Gpos{v+\frac1{2^k}-\frac1{2^{m-1}}}
        {v+\frac1{2^k}}\\
&=\Gpos{\dfrac{v2^m+2^{m-k}-2}{2^m}}
        {\dfrac{v2^m+2^{m-k}}{2^m}}.
\end{aligned}
\]
Theo hàm dyadic,
\[
  G=\frac{v2^m+2^{m-k}-1}{2^m}
   =v+\frac1{2^k}-\frac1{2^{m-1}}+\frac1{2^m}
   =u+\frac1{2^m}.
\]
Vậy thuật toán đúng với mọi \(p\).

Bây giờ xét trạng thái ban đầu \(G\) gồm \(n\) tháp
\(T_1,T_2,\ldots,T_n\), có giá trị tương ứng \(x_1,x_2,\ldots,x_n\). Theo định
lý 8,
\[
  G=x_1+x_2+\cdots+x_n.
\]
Bài toán yêu cầu kiểm tra liệu với mọi tháp \(H\), từ \(C_2+H>0\) có suy ra
\(C_1+H>0\) hay không. Điều này tương đương kiểm tra \(C_1\ge C_2\). Vì vậy
chỉ cần tính tổng giá trị các tháp trong từng cấu hình con rồi so sánh. Độ
phức tạp là \(\Oh(n)\).

\subsubsection*{Tiểu kết}

Ví dụ trên cho thấy surreal number cho một lời giải rất sáng sủa đối với một
lớp trò chơi tổ hợp bất đối xứng có thể tách thành tổng các trò chơi con. Với
\emph{Procrastination}, chương trình theo hướng này rất ngắn. Vì vậy, khi trò
chơi có cấu trúc phù hợp, surreal number nên là lựa chọn đầu tiên.

\section{Quay lại điểm xuất phát}

Surreal number có lợi thế lớn cho một lớp trò chơi bất đối xứng đặc biệt, nhất
là các trò chơi có thể tách thành tổng. Tuy nhiên phạm vi áp dụng của nó khá
hẹp; với những trò chơi tổng quát hơn như các trò chơi cờ thường gặp, công cụ
này không giúp nhiều. Khi đó ta quay lại phân tích trạng thái từ đầu, thiết kế
trạng thái hợp lý và dùng quy hoạch động hoặc lặp.

\subsection{Ví dụ 2: Procrastination}

\subsubsection*{Tóm tắt đề bài}

Đề bài giống ví dụ 1.

\subsubsection*{Phân tích bằng trạng thái}

Từ phần trước, ta biết surreal number giải bài này rất dễ. Nhưng nếu không
dùng công cụ đó, vẫn có thể phân tích trực tiếp.

Định nghĩa W-configuration nói rằng \(L\) thắng dù đi trước hay đi sau. Vì nó
liên quan cả hai khả năng, ta tìm cách đơn giản hóa.

\paragraph{Mệnh đề 3.1.1.}
Với một trạng thái \(G\), nếu \(L\) thắng khi đi trước, thì \(L\) cũng thắng
khi đi sau.

\emph{Chứng minh.}
Nếu \(L\) thắng khi đi trước trong \(G\), tồn tại một lập phương trắng \(K\) mà
nếu \(L\) chọn \(K\), trò chơi chuyển sang trạng thái \(G_1\) trong đó \(R\) đi
trước và thua. Khi \(R\) đi trước trong \(G\), nếu suốt ván \(R\) không động
tới lập phương nào phía trên \(K\), thì \(R\) tương đương đi trước trong
\(G_1\), nên thua. Nếu \(R\) có động tới phần phía trên \(K\), \(L\) lập tức
chọn \(K\) sau nước đó, và \(R\) vẫn tương đương đi trước trong \(G_1\). Vậy
\(R\) thua.

Nhờ kết luận này, chỉ cần xét liệu \(L\) có thắng khi đi trước hay không. Định
nghĩa \(C_1\) không kém \(C_2\) có lượng từ ``mọi tháp \(T\)''. Dù số tháp là
vô hạn, ta thử dùng liệt kê để tìm cấu trúc.

Gọi \(g(X)=1\) nếu \(L\) thắng khi đi trước trong trạng thái \(X\), và
\(g(X)=0\) nếu không. Gọi \(\operatorname{add}(T_1,T_2)\) là tháp tạo bởi việc
đặt \(T_2\) lên trên \(T_1\). Ký hiệu \(T[i]\) là lập phương thứ \(i\) từ dưới
lên; \(W\) là trắng, \(B\) là đen. Gọi mệnh đề \(P\) là ``\(C_1\) không kém
\(C_2\)''.

Cách liệt kê tự nhiên là bắt đầu với tháp rỗng, rồi liên tục đặt thêm một lập
phương trắng hoặc đen lên đỉnh; đây chính là một cây nhị phân sinh các tháp.
Giả sử đang xét một tháp \(T\) với
\[
  g((C_1,T))=g((C_2,T)).
\]
Ta xét việc đặt thêm lập phương lên \(T\).

\paragraph{Mệnh đề 3.1.2.}
Nếu \(g((C,T))=1\), thì với mọi tháp \(T_1\) có \(T_1[1]=W\),
\[
  g((C,\operatorname{add}(T,T_1)))=1.
\]

\emph{Chứng minh.}
Đặt \(T_2=\operatorname{add}(T,T_1)\), chiều cao của \(T\) là \(n\), chiều cao
của \(T_2\) là \(m\). Nếu trong cả ván \(R\) không thao tác lên lập phương nào
phía trên \(T_2[n]\), thì \((C,T_2)\) tương đương \((C,T)\), nên \(L\) thắng khi
đi trước. Nếu \(R\) thao tác phía trên \(T_2[n]\), \(L\) lập tức chọn
\(T_2[n+1]\). Trạng thái ban đầu lại tương đương \((C,T)\), nên \(L\) vẫn
thắng.

Tương tự:

\paragraph{Mệnh đề 3.1.3.}
Nếu \(g((C,T))=0\), thì với mọi tháp \(T_1\) có \(T_1[1]=B\),
\[
  g((C,\operatorname{add}(T,T_1)))=0.
\]

Nếu ở tháp hiện tại \(T\) ta có
\[
  g((C_1,T))=g((C_2,T))=1,
\]
thì nhờ mệnh đề 3.1.2, sau đó chỉ cần liệt kê nhánh đặt thêm một lập phương
đen. Nếu
\[
  g((C_1,T))=g((C_2,T))=0,
\]
thì nhờ mệnh đề 3.1.3, chỉ cần liệt kê nhánh đặt thêm một lập phương trắng.
Nhờ vậy lượng liệt kê giảm từ bậc lũy thừa xuống bậc tuyến tính theo chiều cao
cần xét.

\paragraph{Mệnh đề 3.1.4.}
Nếu \(g((C_1,T))=1\) và \(g((C_2,T))=0\), thì với mọi tháp \(T_1\) thỏa
\[
  g((C_2,\operatorname{add}(T,T_1)))=1,
\]
ta cũng có
\[
  g((C_1,\operatorname{add}(T,T_1)))=1.
\]

\emph{Chứng minh.}
Vì \(g((C_2,T))=0\), nếu
\(g((C_2,\operatorname{add}(T,T_1)))=1\), mệnh đề 3.1.3 cho \(T_1[1]=W\). Lại
có \(g((C_1,T))=1\), nên mệnh đề 3.1.2 cho kết luận.

Do đó, nếu trong quá trình liệt kê gặp \(T\) với
\[
  g((C_1,T))=1,\qquad g((C_2,T))=0,
\]
thì với mọi tháp khiến \((C_2,T)\) là W-configuration, \((C_1,T)\) cũng là
W-configuration; mệnh đề \(P\) đúng. Ngược lại, nếu gặp
\[
  g((C_1,T))=0,\qquad g((C_2,T))=1,
\]
thì \(P\) sai.

Còn phải biết khi nào dừng liệt kê, vì có thể tồn tại \(C_1,C_2\) sao cho với
mọi \(T\),
\[
  g((C_1,T))=g((C_2,T)).
\]
Trường hợp \(C_1=C_2\) là ví dụ đơn giản; hình minh họa trong bản gốc cho thêm
một ví dụ tương tự.

\paragraph{Mệnh đề 3.1.5.}
Giả sử cấu hình con \(C\) gồm ba tháp có chiều cao \(n_1,n_2,n_3\). Khi đó
\(g((C,T))\) chỉ phụ thuộc vào \(n_1+n_2+n_3+1\) lập phương thấp nhất của
\(T\).

Vì mỗi tháp đầu vào có chiều cao không quá \(50\), nếu đã liệt kê tới tháp
\(T\) cao hơn \(151\) mà vẫn có
\[
  g((C_1,T))=g((C_2,T)),
\]
thì với mọi tháp \(T\) hai giá trị đều bằng nhau, và \(P\) đúng.

Tiếp theo là cách hiện thực hàm \(g\). Với trạng thái gồm bốn tháp có chiều
cao \(n_1,n_2,n_3,n_4\), định nghĩa
\[
  d[i,j,k,l,\mathit{cur}]
\]
là trạng thái thắng hay thua khi bốn tháp đã bị giảm xuống các chiều cao
\(i,j,k,l\), và tới lượt người chơi \(\mathit{cur}\). Có thể dùng quy hoạch
động từ dưới lên; với mỗi trạng thái \(S\), liệt kê mọi trạng thái \(S'\) đi
được trong một nước, rồi xác định \(S\) thắng hay thua. Cách trực tiếp có số
trạng thái \(\Oh(n^4)\), mỗi trạng thái tốn \(\Oh(n)\), nên tổng là
\(\Oh(n^5)\), quá lớn.

\paragraph{Mệnh đề 3.1.6.}
Trong một trạng thái \(G\), nếu chọn lập phương \(k\) là một nước thắng cho
người đang đi, thì chọn bất kỳ lập phương cùng màu nào nằm phía trên \(k\)
cũng là một nước thắng.

\emph{Chứng minh.}
Giả sử người đang đi là \(L\), và chọn lập phương trắng \(k\) là nước thắng.
Nếu \(L\) chọn một lập phương trắng \(k'\) phía trên \(k\), thì sau đó nếu
\(R\) không động tới lập phương đen nào phía trên \(k\), ván chơi tương đương
với việc \(L\) chọn \(k\) ngay từ đầu, nên \(L\) thắng. Nếu \(R\) có động tới
phía trên \(k\), \(L\) lập tức chọn \(k\), và ván vẫn tương đương trường hợp
ban đầu. Do đó chọn \(k'\) cũng thắng. Trường hợp của \(R\) đối xứng.

Vì vậy khi tính một trạng thái, ta không cần thử mọi quyết định; trên mỗi tháp
chỉ cần xét lập phương trắng hoặc đen cao nhất phù hợp. Chi phí mỗi trạng thái
giảm xuống \(\Oh(1)\), nên một trạng thái đầy đủ tốn \(\Oh(n^4)\).

Trong quá trình liệt kê tháp mới, tháp mới chỉ khác tháp cũ bởi một lập phương
mới đặt thêm, nên không cần tính lại toàn bộ bảng; chỉ cần tính các trạng thái
mới phát sinh. Vì chiều cao tháp được liệt kê cũng là bậc \(n\), độ phức tạp
toàn bộ là \(\Oh(n^4)\). Với \(n\le 50\), cách này đủ dùng.

\subsubsection*{Tiểu kết}

Cách phân tích này chỉ dùng các trạng thái đặc biệt và quy hoạch động để tính
thắng thua, không cần công cụ quá cao cấp. So với lời giải bằng surreal number
ở ví dụ 1, nó phải phân tích nhiều trường hợp hơn và hiện thực kém gọn hơn.
Tuy nhiên khi gặp một trò chơi mới mà chưa có công cụ sẵn, quay về phân tích
trạng thái cơ bản thường là cách mở đường hiệu quả.

\subsection{Ví dụ 3: The Easy Chase}

\subsubsection*{Tóm tắt đề bài}

Hai người chơi một trò cờ trên bàn \(n\times n\). Người chơi thứ nhất có một
quân trắng ban đầu ở \((w_x,w_y)\); người chơi thứ hai có một quân đen ban đầu
ở \((b_x,b_y)\). \(L\) cầm trắng đi trước, \(R\) cầm đen đi sau.

Khi tới lượt \(L\), quân trắng chọn một trong bốn hướng lên, xuống, trái, phải
và đi đúng một ô. Khi tới lượt \(R\), quân đen cũng chọn một trong bốn hướng
và đi một hoặc hai ô. Không quân nào được đi ra ngoài bàn. Một người thắng khi
nước đi của mình đưa quân của mình tới ô hiện tại của quân đối phương.

Cả hai dùng cùng tiêu chí chiến lược: nếu có thể thắng, họ thắng trong số bước
ít nhất; nếu chắc thua, họ kéo dài ván trong số bước nhiều nhất. Cần dự đoán
người thắng và số bước của ván cờ. Ràng buộc: \(2\le n\le 20\).

\subsubsection*{Phân tích}

Đây là một trò chơi bất đối xứng điển hình; các trò cờ quen thuộc như cờ
tướng hay cờ vua có thể xem là phiên bản phức tạp hơn. Vì bàn cờ nhỏ và chỉ
có hai quân, ta biểu diễn toàn bộ trạng thái có thể xuất hiện rồi tính kết quả
của từng trạng thái.

Một trạng thái được mô tả bằng bộ năm
\[
  (x_1,y_1,x_2,y_2,\mathit{cur}),
\]
trong đó \((x_1,y_1)\) là vị trí quân trắng, \((x_2,y_2)\) là vị trí quân đen,
và \(\mathit{cur}\) cho biết tới lượt \(L\) hay \(R\). Đề bài cần cả người
thắng lẫn số bước, nhưng nếu lưu hai biến \(\mathit{winner}\) và
\(\mathit{move}\) thì chuyển trạng thái hơi cồng kềnh. Ta mã hóa bằng một giá
trị \(f(G)\). Chọn \(\Inf\) là một số nguyên dương rất lớn, chẳng hạn
\(10^8\). Nếu trạng thái \(G\) do \(L\) thắng sau \(x\) bước, đặt
\[
  f(G)=\Inf-x.
\]
Nếu \(G\) do \(R\) thắng sau \(x\) bước, đặt
\[
  f(G)=-\Inf+x.
\]
Ví dụ, nếu \(H\) do \(L\) thắng sau \(5\) bước thì \(f(H)=\Inf-5\). Nếu
\(f(H)=-\Inf+3\), thì \(H\) do \(R\) thắng sau \(3\) bước.

Đặt
\[
  G_L=(w_x,w_y,b_x,b_y,L).
\]
Mục tiêu là tính \(f(G_L)\). Do quan hệ topo giữa các trạng thái không rõ
ràng, ta xác định trước giá trị của các trạng thái kết thúc rồi suy ngược. Có
hai dạng trạng thái kết thúc:
\[
  (x,y,x,y,L),\qquad (x,y,x,y,R).
\]
Khi hai quân đã ở cùng ô và tới lượt \(L\), nghĩa là \(R\) vừa bắt được quân
trắng, nên
\[
  f(x,y,x,y,L)=-\Inf.
\]
Tương tự,
\[
  f(x,y,x,y,R)=\Inf.
\]

Vì \(L\) muốn thắng nhanh nhất nếu thắng và thua muộn nhất nếu thua, mục tiêu
của \(L\) là làm \(f(G)\) lớn nhất. Ngược lại, mục tiêu của \(R\) là làm
\(f(G)\) nhỏ nhất. Ta có các chuyển trạng thái:
\[
\begin{aligned}
f(x_1,y_1,x_2,y_2,L)
&=\max\left\{
  f(x'_1,y'_1,x_2,y_2,R)
  -\sgn\!\left(f(x'_1,y'_1,x_2,y_2,R)\right)
\right\},\\
f(x_1,y_1,x_2,y_2,R)
&=\min\left\{
  f(x_1,y_1,x'_2,y'_2,L)
  -\sgn\!\left(f(x_1,y_1,x'_2,y'_2,L)\right)
\right\}.
\end{aligned}
\]
Ở đây \((x_1,y_1)\ne(x_2,y_2)\); \((x'_1,y'_1)\) là vị trí quân trắng có thể
tới sau một bước; \((x'_2,y'_2)\) là vị trí quân đen có thể tới sau một bước
hợp lệ của nó; và
\[
\sgn(x)=
\begin{cases}
  1, & x>0,\\
  0, & x=0,\\
  -1, & x<0.
\end{cases}
\]

Vấn đề còn lại là thứ tự tính. Vì các trạng thái có thể phụ thuộc lẫn nhau,
khó chọn một thứ tự topo rõ ràng. Tuy nhiên giá trị của một trạng thái không
thể cuối cùng lại suy ngược về chính nó theo kiểu chu trình âm trong bài toán
đường đi ngắn nhất. Vì vậy có thể dùng thuật toán lặp tương tự SPFA:

\begin{verbatim}
Count(G')
    khoi tao f(G) = 0 cho moi trang thai G
    liet ke moi trang thai ket thuc Ge
        xac dinh f(Ge)
        dua Ge vao hang doi q
    while q khong rong
        lay phan tu dau hang doi, goi la Y
        for moi trang thai X co the di mot buoc toi Y
            tmp <- f(X)
            tinh lai f(X) theo cong thuc chuyen trang thai
            if tmp != f(X) va X khong nam trong q
                dua X vao q
    return f(G')
\end{verbatim}

Số trạng thái là \(\Oh(n^4)\), số chuyển cùng bậc với số trạng thái, nên độ
phức tạp là \(\Oh(kn^4)\), trong đó ở trung bình \(k\) là một hằng số nhỏ. Với
\(n\le 20\), thuật toán xử lý tốt bài toán.

\subsubsection*{Mở rộng}

Thuật toán trên không chỉ áp dụng cho bài này. Với nhiều trò chơi có số trạng
thái không lớn, nếu quan hệ giữa các trạng thái rõ ràng, ta thường dùng quy
hoạch động theo thứ tự topo. Khi thứ tự ấy không rõ hoặc khó xác định, lặp
theo kiểu trên là một công cụ mạnh và dễ hiện thực.

\section{Tổng kết}

Muốn giải tốt trò chơi bất đối xứng, cần linh hoạt dùng nhiều cách phân tích
và kết hợp các tư tưởng thuật toán sâu hơn để đạt hiệu quả. Khi thi Tin học
phát triển, các bài toán trò chơi ngày càng đa dạng. Hy vọng các ý tưởng trong
bài viết giúp người đọc có thêm một hướng suy nghĩ khi gặp trò chơi. Để xử lý
tốt trong thời gian thi hạn hẹp, vẫn phải suy nghĩ và thực hành nhiều; kiến
thức trên giấy chỉ thật sự hữu ích khi đã được tự mình rèn luyện.

\paragraph{Lời cảm ơn.}
Tác giả cảm ơn huấn luyện viên Tang Wenbin và Hu Weidong đã hướng dẫn; cảm ơn
thầy Ruan Zhiyuan và Wan Jie vì nhiều năm huấn luyện, hướng dẫn; cảm ơn bạn
Huang Ziqi ở Trường Trung học số 1 Trung Sơn đã hỗ trợ.

\paragraph{Tài liệu tham khảo.}
\begin{enumerate}
  \item J. H. Conway, \emph{On Numbers and Games}.
  \item Claus Tondering, \emph{Surreal Numbers: An Introduction}.
  \item Thomas S. Ferguson, \emph{Game Theory}.
  \item Zhang Yifei, bài viết đội tuyển quốc gia năm 2002 về quá trình giải
  một lớp trò chơi từ trực giác tới lý luận.
  \item Luo Ji, bài viết đội tuyển quốc gia năm 2002 về hai hướng tư duy cho
  bài toán đối sách, bắt đầu từ trò lấy đá.
  \item Wang Xiaoke, bài viết đội tuyển quốc gia năm 2007 về phân tích một lớp
  trò chơi tổ hợp.
\end{enumerate}

\appendix
\section{Đề gốc}

\subsection{Procrastination}

Ngày xưa có hai nghiên cứu sinh là bạn thân. Trong những lúc nghỉ ngắn khỏi
nghiên cứu, thường không quá vài giờ, họ thích chơi trò \emph{Procrastination}.

Đây là trò chơi hai người, đen và trắng, luân phiên đi. Trò chơi thao tác trên
các tháp lập phương; mỗi lập phương có màu đen hoặc trắng. Ban đầu có bốn
tháp, mỗi tháp là một chồng gồm vài lập phương. Ở lượt của mình, người chơi có
thể bỏ bất kỳ lập phương nào cùng màu với mình: người trắng chỉ bỏ lập phương
trắng, người đen chỉ bỏ lập phương đen. Mọi lập phương nằm phía trên lập
phương được chọn cũng bị bỏ khỏi tháp, bất kể màu gì. Ví dụ, nếu một tháp từ
dưới lên là đen, trắng, đen, trắng, thì người đen chọn lập phương đen dưới
cùng sẽ bỏ cả tháp; người đen cũng có thể chọn lập phương thứ ba và bỏ luôn
lập phương thứ tư. Người trắng chọn lập phương thứ hai thì chỉ còn lại một lập
phương đen; người trắng cũng có thể chọn lập phương thứ tư. Người không thể
bỏ lập phương nào sẽ thua.

Hai người đã luyện trò này từ thời đại học và chơi hoàn hảo: nếu một người có
chiến lược thắng, người đó sẽ thắng. Một lúc nào đó, họ phát hiện rằng với hầu
hết cấu hình ban đầu, một người có chiến lược thắng bất kể ai đi trước. Họ gọi
một cấu hình là W-configuration nếu trắng có chiến lược thắng bất kể ai đi
trước, và B-configuration nếu đen có chiến lược thắng bất kể ai đi trước.

Họ cũng nhận thấy một số cấu hình con ít nhất thuận lợi cho một người chơi
bằng những cấu hình con khác. Một cấu hình con \(C\) gồm ba tháp; \(C\) cùng
một tháp thứ tư \(T\) tạo thành một cấu hình đầy đủ, ký hiệu \((C,T)\). Một cấu
hình con \(C_1\) ít nhất thuận lợi cho trắng bằng cấu hình con \(C_2\) nếu và
chỉ nếu với mọi tháp thứ tư \(T\), hễ \((C_2,T)\) là W-configuration thì
\((C_1,T)\) cũng là W-configuration. Nói cách khác, không tồn tại tháp thứ tư
\(T\) sao cho \((C_1,T)\) không phải W-configuration còn \((C_2,T)\) là
W-configuration.

Cho hai cấu hình con \(C_1\) và \(C_2\), hãy kiểm tra liệu \(C_1\) có ít nhất
thuận lợi cho trắng bằng \(C_2\) hay không.

\paragraph{Dữ liệu vào.}
Dòng đầu chứa số bộ kiểm thử. Mỗi bộ gồm một dòng \texttt{Test N}, trong đó
\texttt{N} là số thứ tự bộ kiểm thử, tiếp theo là tám dòng mô tả lần lượt hai
cấu hình con \(C_1,C_2\). Mỗi cấu hình con gồm bốn dòng. Dòng đầu chứa ba số
\(n_1,n_2,n_3\), là chiều cao ba tháp, với
\[
  0\le n_1,n_2,n_3\le 50.
\]
Dòng thứ hai chứa \(n_1\) chữ cái \texttt{B} hoặc \texttt{W}, cách nhau bởi
dấu cách, mô tả tháp thứ nhất. Dòng thứ ba và thứ tư mô tả hai tháp còn lại
tương tự. Chữ \texttt{W} là lập phương trắng, \texttt{B} là lập phương đen.
Mỗi tháp được mô tả từ dưới lên trên.

\paragraph{Dữ liệu ra.}
Với mỗi bộ kiểm thử, in trên một dòng số thứ tự bộ kiểm thử và \texttt{Yes}
nếu \(C_1\) ít nhất thuận lợi cho trắng bằng \(C_2\); ngược lại in
\texttt{No}.

\paragraph{Ví dụ vào.}
\begin{verbatim}
2
Test 1
331
WBB
WBW
B
333
BWW
BWW
WBB
Test 2
332
WBB
WBW
BB
333
BWW
BWW
WBB
\end{verbatim}

\paragraph{Ví dụ ra.}
\begin{verbatim}
Test 1: Yes
Test 2: No
\end{verbatim}

\subsection{The Easy Chase}

Hai người chơi một trò đơn giản trên bàn \(n\times n\). Người chơi thứ nhất có
một quân trắng ban đầu ở \((\mathit{rowWhite},\mathit{colWhite})\). Người chơi
thứ hai có một quân đen ban đầu ở
\((\mathit{rowBlack},\mathit{colBlack})\). Tọa độ đánh số từ \(1\). Hai người
luân phiên đi, người thứ nhất đi trước.

Khi tới lượt người thứ nhất, người đó chọn một trong bốn hướng lên, xuống,
trái, phải và đi quân trắng một ô theo hướng ấy. Khi tới lượt người thứ hai,
người đó cũng chọn một trong bốn hướng và đi quân đen một hoặc hai ô. Một
người thắng khi nước đi của người đó đưa quân của mình tới ô đang có quân đối
phương.

Cả hai dùng chiến lược tối ưu. Nếu có thể thắng, họ chọn chiến lược làm số
nước của ván nhỏ nhất. Nếu không thể thắng, họ chọn chiến lược làm số nước
của ván lớn nhất. Nếu người thứ nhất thắng, trả về \texttt{"WHITE x"}; nếu
người thứ hai thắng, trả về \texttt{"BLACK x"}, trong đó \texttt{x} là số nước
của ván.

\paragraph{Định nghĩa.}
\begin{verbatim}
Class: TheEasyChase
Method: winner
Parameters: int, int, int, int, int
Returns: String
Method signature:
String winner(int n, int rowWhite, int colWhite,
              int rowBlack, int colBlack)
\end{verbatim}
Phương thức phải là \texttt{public}.

\paragraph{Ràng buộc.}
\[
  2\le n\le 20.
\]
Các tọa độ đều nằm trong khoảng từ \(1\) đến \(n\), và hai quân ban đầu nằm ở
hai ô khác nhau.

\paragraph{Ví dụ.}
\begin{verbatim}
0)
n = 2
rowWhite = 1
colWhite = 1
rowBlack = 2
colBlack = 2
Returns: "BLACK 2"
\end{verbatim}
Người thứ nhất có hai nước đi khả dĩ, nhưng kiểu gì cũng thua.

\begin{verbatim}
1)
n = 2
rowWhite = 2
colWhite = 2
rowBlack = 1
colBlack = 2
Returns: "WHITE 1"
\end{verbatim}
Ván này kết thúc chỉ sau một nước.

\begin{verbatim}
2)
n = 3
rowWhite = 1
colWhite = 1
rowBlack = 3
colBlack = 3
Returns: "BLACK 6"
\end{verbatim}

\section{Chương trình tham khảo}

Trong lớp văn bản trích xuất được từ PDF bằng \texttt{pdftotext}, phần phụ lục
chương trình chỉ còn các nhan đề ``chương trình ví dụ 1'', ``chương trình ví
dụ 2'' và ``chương trình ví dụ 3'', không có thân mã nguồn. Vì vậy bản dịch
giữ phần phụ lục ở mức ghi chú này và không dựng lại chương trình từ suy đoán.

\end{document}
